% Options for packages loaded elsewhere
\PassOptionsToPackage{unicode}{hyperref}
\PassOptionsToPackage{hyphens}{url}
\documentclass[
]{article}
\usepackage{xcolor}
\usepackage[margin=1in]{geometry}
\usepackage{amsmath,amssymb}
\setcounter{secnumdepth}{-\maxdimen} % remove section numbering
\usepackage{iftex}
\ifPDFTeX
  \usepackage[T1]{fontenc}
  \usepackage[utf8]{inputenc}
  \usepackage{textcomp} % provide euro and other symbols
\else % if luatex or xetex
  \usepackage{unicode-math} % this also loads fontspec
  \defaultfontfeatures{Scale=MatchLowercase}
  \defaultfontfeatures[\rmfamily]{Ligatures=TeX,Scale=1}
\fi
\usepackage{lmodern}
\ifPDFTeX\else
  % xetex/luatex font selection
\fi
% Use upquote if available, for straight quotes in verbatim environments
\IfFileExists{upquote.sty}{\usepackage{upquote}}{}
\IfFileExists{microtype.sty}{% use microtype if available
  \usepackage[]{microtype}
  \UseMicrotypeSet[protrusion]{basicmath} % disable protrusion for tt fonts
}{}
\makeatletter
\@ifundefined{KOMAClassName}{% if non-KOMA class
  \IfFileExists{parskip.sty}{%
    \usepackage{parskip}
  }{% else
    \setlength{\parindent}{0pt}
    \setlength{\parskip}{6pt plus 2pt minus 1pt}}
}{% if KOMA class
  \KOMAoptions{parskip=half}}
\makeatother
\usepackage{color}
\usepackage{fancyvrb}
\newcommand{\VerbBar}{|}
\newcommand{\VERB}{\Verb[commandchars=\\\{\}]}
\DefineVerbatimEnvironment{Highlighting}{Verbatim}{commandchars=\\\{\}}
% Add ',fontsize=\small' for more characters per line
\usepackage{framed}
\definecolor{shadecolor}{RGB}{248,248,248}
\newenvironment{Shaded}{\begin{snugshade}}{\end{snugshade}}
\newcommand{\AlertTok}[1]{\textcolor[rgb]{0.94,0.16,0.16}{#1}}
\newcommand{\AnnotationTok}[1]{\textcolor[rgb]{0.56,0.35,0.01}{\textbf{\textit{#1}}}}
\newcommand{\AttributeTok}[1]{\textcolor[rgb]{0.13,0.29,0.53}{#1}}
\newcommand{\BaseNTok}[1]{\textcolor[rgb]{0.00,0.00,0.81}{#1}}
\newcommand{\BuiltInTok}[1]{#1}
\newcommand{\CharTok}[1]{\textcolor[rgb]{0.31,0.60,0.02}{#1}}
\newcommand{\CommentTok}[1]{\textcolor[rgb]{0.56,0.35,0.01}{\textit{#1}}}
\newcommand{\CommentVarTok}[1]{\textcolor[rgb]{0.56,0.35,0.01}{\textbf{\textit{#1}}}}
\newcommand{\ConstantTok}[1]{\textcolor[rgb]{0.56,0.35,0.01}{#1}}
\newcommand{\ControlFlowTok}[1]{\textcolor[rgb]{0.13,0.29,0.53}{\textbf{#1}}}
\newcommand{\DataTypeTok}[1]{\textcolor[rgb]{0.13,0.29,0.53}{#1}}
\newcommand{\DecValTok}[1]{\textcolor[rgb]{0.00,0.00,0.81}{#1}}
\newcommand{\DocumentationTok}[1]{\textcolor[rgb]{0.56,0.35,0.01}{\textbf{\textit{#1}}}}
\newcommand{\ErrorTok}[1]{\textcolor[rgb]{0.64,0.00,0.00}{\textbf{#1}}}
\newcommand{\ExtensionTok}[1]{#1}
\newcommand{\FloatTok}[1]{\textcolor[rgb]{0.00,0.00,0.81}{#1}}
\newcommand{\FunctionTok}[1]{\textcolor[rgb]{0.13,0.29,0.53}{\textbf{#1}}}
\newcommand{\ImportTok}[1]{#1}
\newcommand{\InformationTok}[1]{\textcolor[rgb]{0.56,0.35,0.01}{\textbf{\textit{#1}}}}
\newcommand{\KeywordTok}[1]{\textcolor[rgb]{0.13,0.29,0.53}{\textbf{#1}}}
\newcommand{\NormalTok}[1]{#1}
\newcommand{\OperatorTok}[1]{\textcolor[rgb]{0.81,0.36,0.00}{\textbf{#1}}}
\newcommand{\OtherTok}[1]{\textcolor[rgb]{0.56,0.35,0.01}{#1}}
\newcommand{\PreprocessorTok}[1]{\textcolor[rgb]{0.56,0.35,0.01}{\textit{#1}}}
\newcommand{\RegionMarkerTok}[1]{#1}
\newcommand{\SpecialCharTok}[1]{\textcolor[rgb]{0.81,0.36,0.00}{\textbf{#1}}}
\newcommand{\SpecialStringTok}[1]{\textcolor[rgb]{0.31,0.60,0.02}{#1}}
\newcommand{\StringTok}[1]{\textcolor[rgb]{0.31,0.60,0.02}{#1}}
\newcommand{\VariableTok}[1]{\textcolor[rgb]{0.00,0.00,0.00}{#1}}
\newcommand{\VerbatimStringTok}[1]{\textcolor[rgb]{0.31,0.60,0.02}{#1}}
\newcommand{\WarningTok}[1]{\textcolor[rgb]{0.56,0.35,0.01}{\textbf{\textit{#1}}}}
\usepackage{longtable,booktabs,array}
\newcounter{none} % for unnumbered tables
\usepackage{calc} % for calculating minipage widths
% Correct order of tables after \paragraph or \subparagraph
\usepackage{etoolbox}
\makeatletter
\patchcmd\longtable{\par}{\if@noskipsec\mbox{}\fi\par}{}{}
\makeatother
% Allow footnotes in longtable head/foot
\IfFileExists{footnotehyper.sty}{\usepackage{footnotehyper}}{\usepackage{footnote}}
\makesavenoteenv{longtable}
\usepackage{graphicx}
\makeatletter
\newsavebox\pandoc@box
\newcommand*\pandocbounded[1]{% scales image to fit in text height/width
  \sbox\pandoc@box{#1}%
  \Gscale@div\@tempa{\textheight}{\dimexpr\ht\pandoc@box+\dp\pandoc@box\relax}%
  \Gscale@div\@tempb{\linewidth}{\wd\pandoc@box}%
  \ifdim\@tempb\p@<\@tempa\p@\let\@tempa\@tempb\fi% select the smaller of both
  \ifdim\@tempa\p@<\p@\scalebox{\@tempa}{\usebox\pandoc@box}%
  \else\usebox{\pandoc@box}%
  \fi%
}
% Set default figure placement to htbp
\def\fps@figure{htbp}
\makeatother
\setlength{\emergencystretch}{3em} % prevent overfull lines
\providecommand{\tightlist}{%
  \setlength{\itemsep}{0pt}\setlength{\parskip}{0pt}}
\usepackage{bookmark}
\IfFileExists{xurl.sty}{\usepackage{xurl}}{} % add URL line breaks if available
\urlstyle{same}
\hypersetup{
  pdftitle={📊 Aventura Estatística na Biblioteca Omar Catunda},
  pdfauthor={Seu nome},
  hidelinks,
  pdfcreator={LaTeX via pandoc}}

\title{📊 Aventura Estatística na Biblioteca Omar Catunda}
\author{Seu nome}
\date{2026-09-23}

\begin{document}
\maketitle

{
\setcounter{tocdepth}{2}
\tableofcontents
}
\begin{quote}
\textbf{🎯 Missão}

Vocês vão até a \textbf{Biblioteca Universitária de Ciências e
Tecnologias Professor Omar Catunda} (UFBA, Campus de Ondina) para
investigar uma pergunta curiosa:

\textbf{Os livros de Estatística são todos ``iguais''? Ou será que até
dentro da mesma área existem diferenças?}

A atividade acontece \textbf{no horário da aula (10:40--12:30)}. A
biblioteca está aberta de segunda a sexta, das 7h às 19h, então vocês
têm tempo de sobra.

O objetivo é duplo: 1. Aplicar na prática um \textbf{Delineamento
Inteiramente Casualizado (DIC)}. 2. \textbf{Conhecer o acervo de
Estatística} da biblioteca e descobrir o que ele tem de interessante.
\end{quote}

\begin{center}\rule{0.5\linewidth}{0.5pt}\end{center}

\subsection{🗺️ Parte 1 --- Localizando o
tesouro}\label{parte-1-localizando-o-tesouro}

A biblioteca fica no \textbf{Campus de Ondina, em frente ao Instituto de
Física} e atrás da Faculdade de Arquitetura. O acervo de Estatística
está no \textbf{2º andar}, junto com as demais áreas de exatas
(Matemática, Computação, Física, Química, Geociências).

\subsubsection{Informações gerais}\label{informauxe7uxf5es-gerais}

{\def\LTcaptype{none} % do not increment counter
\begin{longtable}[]{@{}
  >{\raggedright\arraybackslash}p{(\linewidth - 2\tabcolsep) * \real{0.3333}}
  >{\raggedright\arraybackslash}p{(\linewidth - 2\tabcolsep) * \real{0.6667}}@{}}
\toprule\noalign{}
\begin{minipage}[b]{\linewidth}\raggedright
Item
\end{minipage} & \begin{minipage}[b]{\linewidth}\raggedright
Informação
\end{minipage} \\
\midrule\noalign{}
\endhead
\bottomrule\noalign{}
\endlastfoot
Endereço & Rua Barão de Jeremoabo, s/n --- Campus de Ondina \\
Horário & Segunda a sexta, \textbf{7h às 19h} \\
Andar do acervo & \textbf{2º andar} \\
O que vocês vão encontrar & Livros de Estatística, Matemática, Física,
Química, Computação \\
\end{longtable}
}

\subsubsection{Tarefa rápida de exploração (5
minutos)}\label{tarefa-ruxe1pida-de-explorauxe7uxe3o-5-minutos}

Antes de começar a medir, localizem a \textbf{seção de Estatística}.
Vocês devem encontrar \textbf{3 prateleiras} (ou 3 seções de estante)
com livros dessa área.

\begin{enumerate}
\def\labelenumi{\arabic{enumi}.}
\tightlist
\item
  \textbf{As 3 prateleiras de Estatística estão próximas ou separadas?}
\item
  \textbf{Vocês reconhecem algum livro?} (Bussab \& Morettin,
  Montgomery, etc.)
\item
  \textbf{Qual prateleira parece ter os livros mais antigos?}
\end{enumerate}

\begin{quote}
💡 \textbf{Importante}: o ``tratamento'' do experimento serão as
\textbf{3 prateleiras de Estatística}. Cada prateleira é um nível
diferente.
\end{quote}

\begin{center}\rule{0.5\linewidth}{0.5pt}\end{center}

\subsection{🧪 Parte 2 --- O experimento
(DIC)}\label{parte-2-o-experimento-dic}

\subsubsection{Hipóteses}\label{hipuxf3teses}

\begin{quote}
\textbf{H₀ (hipótese nula)}: a média de páginas dos livros de
Estatística é igual entre as 3 prateleiras.

\textbf{H₁ (hipótese alternativa)}: pelo menos uma prateleira tem média
de páginas diferente das demais.
\end{quote}

\subsubsection{Estrutura do
delineamento}\label{estrutura-do-delineamento}

{\def\LTcaptype{none} % do not increment counter
\begin{longtable}[]{@{}
  >{\raggedright\arraybackslash}p{(\linewidth - 2\tabcolsep) * \real{0.4762}}
  >{\raggedright\arraybackslash}p{(\linewidth - 2\tabcolsep) * \real{0.5238}}@{}}
\toprule\noalign{}
\begin{minipage}[b]{\linewidth}\raggedright
Elemento
\end{minipage} & \begin{minipage}[b]{\linewidth}\raggedright
Definição
\end{minipage} \\
\midrule\noalign{}
\endhead
\bottomrule\noalign{}
\endlastfoot
\textbf{Unidade experimental} & Um livro de Estatística \\
\textbf{Tratamento} & Prateleira de Estatística (3 níveis: P1, P2,
P3) \\
\textbf{Variável resposta} & Número de páginas do livro \\
\textbf{Replicações por tratamento} & \textbf{8 livros} por prateleira
(total = 24 livros) \\
\textbf{Total de alunos} & 8 pessoas → \textbf{cada aluno mede 3
livros} \\
\end{longtable}
}

\begin{quote}
📌 \textbf{Por que ``número de páginas''?} Porque é rápido de medir, não
precisa de instrumento, não danifica o livro, e revela diferenças
interessantes entre tipos de livro (manuais densos vs.~textos
introdutórios).
\end{quote}

\subsubsection{Plano de
aleatorização}\label{plano-de-aleatorizauxe7uxe3o}

\begin{Shaded}
\begin{Highlighting}[]
\CommentTok{\# Cada grupo usa uma semente diferente. Sugestão: use o número da matrícula.}
\FunctionTok{set.seed}\NormalTok{(}\DecValTok{2025}\NormalTok{)  }\CommentTok{\# ← troquem pelo número da sorte do grupo}

\CommentTok{\# As 3 prateleiras de Estatística (identifiquem fisicamente: P1, P2, P3)}
\NormalTok{prateleiras }\OtherTok{\textless{}{-}} \FunctionTok{c}\NormalTok{(}\StringTok{"Prateleira 1"}\NormalTok{, }\StringTok{"Prateleira 2"}\NormalTok{, }\StringTok{"Prateleira 3"}\NormalTok{)}

\CommentTok{\# Quantos livros medir por prateleira}
\NormalTok{n\_por\_prateleira }\OtherTok{\textless{}{-}} \DecValTok{8}

\NormalTok{plano }\OtherTok{\textless{}{-}} \FunctionTok{expand.grid}\NormalTok{(}
  \AttributeTok{prateleira =}\NormalTok{ prateleiras, }
  \AttributeTok{livro\_id =} \DecValTok{1}\SpecialCharTok{:}\NormalTok{n\_por\_prateleira}
\NormalTok{) }\SpecialCharTok{|\textgreater{}}
  \FunctionTok{mutate}\NormalTok{(}
    \CommentTok{\# Ordem aleatória de medição (evita viés de cansaço)}
    \AttributeTok{ordem\_medicao =} \FunctionTok{sample}\NormalTok{(}\DecValTok{1}\SpecialCharTok{:}\FunctionTok{n}\NormalTok{()),}
    \CommentTok{\# Posição aproximada na prateleira (para orientar a escolha)}
    \AttributeTok{posicao =} \FunctionTok{sample}\NormalTok{(}\DecValTok{1}\SpecialCharTok{:}\DecValTok{15}\NormalTok{, }\FunctionTok{n}\NormalTok{(), }\AttributeTok{replace =} \ConstantTok{TRUE}\NormalTok{),}
    \CommentTok{\# Aluno responsável (distribuam entre vocês)}
    \AttributeTok{aluno =} \FunctionTok{rep}\NormalTok{(}\DecValTok{1}\SpecialCharTok{:}\DecValTok{8}\NormalTok{, }\AttributeTok{length.out =} \FunctionTok{n}\NormalTok{())}
\NormalTok{  ) }\SpecialCharTok{|\textgreater{}}
  \FunctionTok{arrange}\NormalTok{(ordem\_medicao)}

\FunctionTok{kable}\NormalTok{(plano, }\AttributeTok{caption =} \StringTok{"📋 Plano de amostragem e ordem de medição"}\NormalTok{)}
\end{Highlighting}
\end{Shaded}

\begin{longtable}[]{@{}lrrrr@{}}
\caption{📋 Plano de amostragem e ordem de medição}\tabularnewline
\toprule\noalign{}
prateleira & livro\_id & ordem\_medicao & posicao & aluno \\
\midrule\noalign{}
\endfirsthead
\toprule\noalign{}
prateleira & livro\_id & ordem\_medicao & posicao & aluno \\
\midrule\noalign{}
\endhead
\bottomrule\noalign{}
\endlastfoot
Prateleira 1 & 2 & 1 & 4 & 4 \\
Prateleira 3 & 7 & 2 & 13 & 5 \\
Prateleira 2 & 7 & 3 & 10 & 4 \\
Prateleira 3 & 1 & 4 & 5 & 3 \\
Prateleira 1 & 5 & 5 & 11 & 5 \\
Prateleira 1 & 7 & 6 & 4 & 3 \\
Prateleira 2 & 4 & 7 & 2 & 3 \\
Prateleira 2 & 8 & 8 & 7 & 7 \\
Prateleira 3 & 6 & 9 & 8 & 2 \\
Prateleira 2 & 2 & 10 & 10 & 5 \\
Prateleira 1 & 4 & 11 & 2 & 2 \\
Prateleira 2 & 1 & 12 & 3 & 2 \\
Prateleira 1 & 1 & 13 & 13 & 1 \\
Prateleira 2 & 5 & 14 & 11 & 6 \\
Prateleira 3 & 4 & 15 & 9 & 4 \\
Prateleira 2 & 6 & 16 & 4 & 1 \\
Prateleira 1 & 3 & 17 & 1 & 7 \\
Prateleira 3 & 5 & 18 & 1 & 7 \\
Prateleira 3 & 8 & 19 & 3 & 8 \\
Prateleira 1 & 8 & 20 & 11 & 6 \\
Prateleira 1 & 6 & 21 & 14 & 8 \\
Prateleira 3 & 3 & 22 & 15 & 1 \\
Prateleira 2 & 3 & 23 & 3 & 8 \\
Prateleira 3 & 2 & 24 & 2 & 6 \\
\end{longtable}

\begin{quote}
\textbf{Como usar:} - \textbf{Antes da aula}: rodem o código e
imprimam/salvem a tabela. - \textbf{Na biblioteca}: sigam a coluna
\texttt{ordem\_medicao}. - Para cada linha, vão até a prateleira
indicada, escolham \textbf{um livro qualquer} na posição aproximada (não
escolham pelo tamanho!), e anotem o \textbf{número de páginas}.
\end{quote}

\begin{center}\rule{0.5\linewidth}{0.5pt}\end{center}

\subsection{📝 Parte 3 --- Coleta de dados (em sala,
10:40--12:30)}\label{parte-3-coleta-de-dados-em-sala-10401230}

Cada aluno é responsável por \textbf{3 livros} (total 24). Preencham a
tabela com os dados coletados.

\begin{Shaded}
\begin{Highlighting}[]
\CommentTok{\# Preencham com os dados reais após a visita}
\NormalTok{dados }\OtherTok{\textless{}{-}} \FunctionTok{tibble}\NormalTok{(}
  \AttributeTok{prateleira =} \FunctionTok{c}\NormalTok{(}
    \FunctionTok{rep}\NormalTok{(}\StringTok{"Prateleira 1"}\NormalTok{, }\DecValTok{8}\NormalTok{),}
    \FunctionTok{rep}\NormalTok{(}\StringTok{"Prateleira 2"}\NormalTok{, }\DecValTok{8}\NormalTok{),}
    \FunctionTok{rep}\NormalTok{(}\StringTok{"Prateleira 3"}\NormalTok{, }\DecValTok{8}\NormalTok{)}
\NormalTok{  ),}
  \AttributeTok{livro\_id =} \FunctionTok{rep}\NormalTok{(}\DecValTok{1}\SpecialCharTok{:}\DecValTok{8}\NormalTok{, }\DecValTok{3}\NormalTok{),}
  \AttributeTok{paginas =} \FunctionTok{c}\NormalTok{(}
    \CommentTok{\# 👇 COLOQUEM AQUI os números de páginas medidos}
    \ConstantTok{NA}\NormalTok{, }\ConstantTok{NA}\NormalTok{, }\ConstantTok{NA}\NormalTok{, }\ConstantTok{NA}\NormalTok{, }\ConstantTok{NA}\NormalTok{, }\ConstantTok{NA}\NormalTok{, }\ConstantTok{NA}\NormalTok{, }\ConstantTok{NA}\NormalTok{,   }\CommentTok{\# Prateleira 1}
    \ConstantTok{NA}\NormalTok{, }\ConstantTok{NA}\NormalTok{, }\ConstantTok{NA}\NormalTok{, }\ConstantTok{NA}\NormalTok{, }\ConstantTok{NA}\NormalTok{, }\ConstantTok{NA}\NormalTok{, }\ConstantTok{NA}\NormalTok{, }\ConstantTok{NA}\NormalTok{,   }\CommentTok{\# Prateleira 2}
    \ConstantTok{NA}\NormalTok{, }\ConstantTok{NA}\NormalTok{, }\ConstantTok{NA}\NormalTok{, }\ConstantTok{NA}\NormalTok{, }\ConstantTok{NA}\NormalTok{, }\ConstantTok{NA}\NormalTok{, }\ConstantTok{NA}\NormalTok{, }\ConstantTok{NA}    \CommentTok{\# Prateleira 3}
\NormalTok{  )}
\NormalTok{)}
\end{Highlighting}
\end{Shaded}

\begin{quote}
\textbf{⏱️ Cronograma sugerido para 10:40--12:30:} \textbar{} Tempo
\textbar{} Atividade \textbar{}
\textbar-------\textbar-----------\textbar{} \textbar{} 10:40--10:50
\textbar{} localização das prateleiras \textbar{} \textbar{}
10:50--11:10 \textbar{} Coleta de dados (medição de páginas) \textbar{}
\textbar{} 11:10--11:40 \textbar{} Digitação e organização dos dados
\textbar{} \textbar{} 11:40--12:10 \textbar{} Análise no R (descritiva,
ANOVA, gráficos) \textbar{} \textbar{} 12:10--12:30 \textbar{} Discussão
e conclusões \textbar{}
\end{quote}

\begin{center}\rule{0.5\linewidth}{0.5pt}\end{center}

\subsection{📊 Parte 4 --- Análise dos
dados}\label{parte-4-anuxe1lise-dos-dados}

\subsubsection{Estatísticas
descritivas}\label{estatuxedsticas-descritivas}

\begin{Shaded}
\begin{Highlighting}[]
\NormalTok{dados }\SpecialCharTok{|\textgreater{}}
  \FunctionTok{group\_by}\NormalTok{(prateleira) }\SpecialCharTok{|\textgreater{}}
  \FunctionTok{summarise}\NormalTok{(}
    \AttributeTok{n =} \FunctionTok{n}\NormalTok{(),}
    \AttributeTok{media =} \FunctionTok{mean}\NormalTok{(paginas, }\AttributeTok{na.rm =} \ConstantTok{TRUE}\NormalTok{),}
    \AttributeTok{desvio =} \FunctionTok{sd}\NormalTok{(paginas, }\AttributeTok{na.rm =} \ConstantTok{TRUE}\NormalTok{),}
    \AttributeTok{mediana =} \FunctionTok{median}\NormalTok{(paginas, }\AttributeTok{na.rm =} \ConstantTok{TRUE}\NormalTok{),}
    \AttributeTok{minimo =} \FunctionTok{min}\NormalTok{(paginas, }\AttributeTok{na.rm =} \ConstantTok{TRUE}\NormalTok{),}
    \AttributeTok{maximo =} \FunctionTok{max}\NormalTok{(paginas, }\AttributeTok{na.rm =} \ConstantTok{TRUE}\NormalTok{)}
\NormalTok{  ) }\SpecialCharTok{|\textgreater{}}
  \FunctionTok{kable}\NormalTok{(}\AttributeTok{caption =} \StringTok{"Estatísticas descritivas por prateleira"}\NormalTok{)}
\end{Highlighting}
\end{Shaded}

\subsubsection{Gráfico}\label{gruxe1fico}

\begin{Shaded}
\begin{Highlighting}[]
\FunctionTok{ggplot}\NormalTok{(dados, }\FunctionTok{aes}\NormalTok{(}\AttributeTok{x =}\NormalTok{ prateleira, }\AttributeTok{y =}\NormalTok{ paginas, }\AttributeTok{fill =}\NormalTok{ prateleira)) }\SpecialCharTok{+}
  \FunctionTok{geom\_boxplot}\NormalTok{(}\AttributeTok{alpha =} \FloatTok{0.7}\NormalTok{, }\AttributeTok{show.legend =} \ConstantTok{FALSE}\NormalTok{) }\SpecialCharTok{+}
  \FunctionTok{geom\_jitter}\NormalTok{(}\AttributeTok{width =} \FloatTok{0.15}\NormalTok{, }\AttributeTok{size =} \DecValTok{2}\NormalTok{, }\AttributeTok{alpha =} \FloatTok{0.5}\NormalTok{) }\SpecialCharTok{+}
  \FunctionTok{labs}\NormalTok{(}
    \AttributeTok{title =} \StringTok{"📚 Distribuição do número de páginas por prateleira"}\NormalTok{,}
    \AttributeTok{x =} \StringTok{"Prateleira de Estatística"}\NormalTok{,}
    \AttributeTok{y =} \StringTok{"Número de páginas"}
\NormalTok{  ) }\SpecialCharTok{+}
  \FunctionTok{theme\_minimal}\NormalTok{(}\AttributeTok{base\_size =} \DecValTok{14}\NormalTok{) }\SpecialCharTok{+}
  \FunctionTok{scale\_fill\_brewer}\NormalTok{(}\AttributeTok{palette =} \StringTok{"Set2"}\NormalTok{)}
\end{Highlighting}
\end{Shaded}

\subsubsection{ANOVA (um fator)}\label{anova-um-fator}

\begin{Shaded}
\begin{Highlighting}[]
\NormalTok{modelo }\OtherTok{\textless{}{-}} \FunctionTok{aov}\NormalTok{(paginas }\SpecialCharTok{\textasciitilde{}}\NormalTok{ prateleira, }\AttributeTok{data =}\NormalTok{ dados)}
\FunctionTok{summary}\NormalTok{(modelo)}
\end{Highlighting}
\end{Shaded}

\subsubsection{Comparações múltiplas (se
significativo)}\label{comparauxe7uxf5es-muxfaltiplas-se-significativo}

\begin{Shaded}
\begin{Highlighting}[]
\FunctionTok{TukeyHSD}\NormalTok{(modelo)}
\end{Highlighting}
\end{Shaded}

\subsubsection{Verificação dos
pressupostos}\label{verificauxe7uxe3o-dos-pressupostos}

\begin{Shaded}
\begin{Highlighting}[]
\CommentTok{\# Normalidade dos resíduos}
\FunctionTok{shapiro.test}\NormalTok{(}\FunctionTok{residuals}\NormalTok{(modelo))}

\CommentTok{\# Homocedasticidade}
\FunctionTok{library}\NormalTok{(car)}
\FunctionTok{leveneTest}\NormalTok{(paginas }\SpecialCharTok{\textasciitilde{}}\NormalTok{ prateleira, }\AttributeTok{data =}\NormalTok{ dados)}
\end{Highlighting}
\end{Shaded}

\begin{center}\rule{0.5\linewidth}{0.5pt}\end{center}

\subsection{🎮 Parte 5 --- Desafios extras (se sobrar
tempo)}\label{parte-5-desafios-extras-se-sobrar-tempo}

{\def\LTcaptype{none} % do not increment counter
\begin{longtable}[]{@{}
  >{\raggedright\arraybackslash}p{(\linewidth - 2\tabcolsep) * \real{0.4091}}
  >{\raggedright\arraybackslash}p{(\linewidth - 2\tabcolsep) * \real{0.5909}}@{}}
\toprule\noalign{}
\begin{minipage}[b]{\linewidth}\raggedright
Desafio
\end{minipage} & \begin{minipage}[b]{\linewidth}\raggedright
O que fazer
\end{minipage} \\
\midrule\noalign{}
\endhead
\bottomrule\noalign{}
\endlastfoot
🕵️ \textbf{Caçadores de livros raros} & Qual é o livro de Estatística
mais antigo que vocês encontraram? \\
📏 \textbf{O mais grosso e o mais fino} & Quais são os extremos de
espessura? Em qual prateleira estão? \\
🤔 \textbf{Meta-experimento} & Se refizessem o experimento, o que
mudariam? Medir espessura? Incluir mais prateleiras? \\
\end{longtable}
}

\begin{center}\rule{0.5\linewidth}{0.5pt}\end{center}

\subsection{📤 O que entregar}\label{o-que-entregar}

Cada grupo entrega \textbf{1 relatório em HTML}, contendo:

\begin{enumerate}
\def\labelenumi{\arabic{enumi}.}
\tightlist
\item
  Nomes das 3 prateleiras escolhidas (com descrição de onde estão).
\item
  A tabela \texttt{plano} completa (aleatorização).
\item
  A tabela \texttt{dados} preenchida.
\item
  Tabela descritiva + boxplot.
\item
  Resultado da ANOVA (valor de p).
\item
  \textbf{Um parágrafo de conclusão}: houve diferença significativa
  entre as prateleiras? O que isso significa? Quais as limitações?
\end{enumerate}

\begin{quote}
\textbf{Dica final}: diferenças no número de páginas podem refletir
\textbf{tipos diferentes de livro} --- manuais mais densos (ex.:
Montgomery) vs.~textos introdutórios (ex.: Bussab \& Morettin), ou
edições mais antigas vs.~recentes. Isso não é ``erro'' do experimento
--- é exatamente o tipo de discussão que vale a pena em Planejamento de
Experimentos.
\end{quote}

\begin{center}\rule{0.5\linewidth}{0.5pt}\end{center}

\subsection{📚 Anexo --- Curiosidades sobre a
biblioteca}\label{anexo-curiosidades-sobre-a-biblioteca}

\begin{itemize}
\tightlist
\item
  A biblioteca leva o nome do matemático \textbf{Omar Catunda}, que veio
  para Salvador em 1960 para dirigir o Instituto de Matemática e Física
  da UFBA.
\item
  Foi inaugurada em \textbf{2 de agosto de 2016}, integrando as
  comemorações dos 70 anos da UFBA.
\item
  O prédio tem \textbf{4,6 mil m²} e reúne cerca de \textbf{40 mil
  volumes} nas áreas de exatas.
\item
  O acervo de Estatística está no \textbf{2º andar}, junto com
  Matemática, Computação, Física, Química e Geociências.
\end{itemize}

\begin{center}\rule{0.5\linewidth}{0.5pt}\end{center}

\emph{Boa aventura estatística! Lembrem-se: um bom experimento não é
aquele que ``dá certo'', é aquele que faz vocês entenderem o porquê.}
```

\begin{center}\rule{0.5\linewidth}{0.5pt}\end{center}

\end{document}
